% !TeX TXS-program:compile = txs:///arara
% arara: lualatex: {shell: yes, synctex: no, interaction: batchmode}
% arara: lualatex: {shell: yes, synctex: no, interaction: batchmode}
% arara: lualatex: {shell: yes, synctex: no, interaction: batchmode}

\documentclass[a4paper,english,11pt]{article}
\def\PLversion{4.00d}
\def\PLdate{15/02/2026}
\usepackage{amsfonts}
\usepackage[nonamssymb,warningsoff,nosiunitxfr]{ProfLycee}
\useproflyclib{piton,ecritures}
\usepackage[executable=python,ignoreerrors]{pyluatex}
\usepackage[math-style=french]{fourier-otf}
\usepackage{mathrsfs}%pour mathscr
\usepackage{awesomebox}
\usepackage[lua]{tkz-euclide}
\usepackage{tkz-tab}
\tikzstyle{every picture}+=[remember picture]
\usetikzlibrary{hobby}
\usepackage{siunitx}
\sisetup{locale=US}
\usepackage{enumitem}
\usepackage{openmoji}
\usepackage{ulem}
\usepackage{fancyvrb}
\usepackage{fancyhdr}
\usepackage{tabularray}
\usepackage{multicol}
\DeclareMathSymbol{;}\mathbin{operators}{'73} %espacement avec ;
%fancy
\fancyhf{}
\renewcommand{\headrulewidth}{0pt}
\lfoot{\sffamily \small [ProfLycee]}
\cfoot{\sffamily \small - \thepage{} -}
\rfoot{\hyperlink{matoc}{\small\faArrowAltCircleUp[regular]}}

\usepackage{codehigh}
\usepackage{graphics}
\usepackage{hologo}
\providecommand\tikzlogo{Ti\textit{k}Z}
\providecommand\TeXLive{\TeX{}Live\xspace}
\providecommand\PSTricks{\textsf{PSTricks}\xspace}
\let\pstricks\PSTricks
\let\TikZ\tikzlogo
\usepackage{menukeys}
\let\tab\relax
\usepackage{tabto}
\usepackage{pgf,pgfplots}
\pgfplotsset{compat=newest,xlabel near ticks,ylabel near ticks}
\usepackage{scontents}
\usepackage{geometry}
\geometry{margin=2cm}
\setlength\parindent{0pt}
\usepackage{babel}
\usepackage{newverbs}

\date{\PLdate}
\newcommand*{\version}{\PLversion}
\title{ProfLycee (\textit{simple doc}) \\User’s Guide for \LaTeX\ (v\version)}
\author{Cédric Pierquet \\ \texttt{cpierquet@outlook.fr}}

%present macro
\tcbuselibrary{minted}
\newtcblisting{myprestexcode}[2]{width=0.8\linewidth,colback=violet!1!white,colframe=violet!75!white,fonttitle=\footnotesize\bfseries\sffamily,title={\vphantom{(qÉ)}\enflag\ #1 | \frflag\ #2\hfill{}[\textit{Syntax}]},listing only,left=1mm,right=1mm,top=1mm,bottom=1mm,minted options={fontsize=\footnotesize,autogobble,tabsize=2}}
\newtcblisting{myprestex}[2]{width=0.8\linewidth,colback=teal!1!white,colframe=teal!75!white,fonttitle=\footnotesize\bfseries\sffamily,title={\vphantom{(qÉ)}\enflag\ #1 | \frflag\ #2\hfill{}[\textit{Sample(s)}]},left=1mm,right=1mm,top=1mm,bottom=1mm,fontupper=\footnotesize,minted options={fontsize=\footnotesize,autogobble,tabsize=2}}

\newcommand\frflag{\openmoji[dstrut=(qH)]{flag- france}}
\newcommand\enflag{\openmoji[dstrut=(qH)]{flag- united states}}

\usepackage{hyperref}
\urlstyle{same}
\hypersetup{pdfborder=0 0 0}

\begin{document}
\maketitle

\bigskip

The purpose of this simplified documentation is to present the main macros of \texttt{ProfLycee}, including:

\begin{itemize}
  \item available names (base / French alias / English alias) if available;
  \item available keys (French / English);
  \item arguments (optional or not).
\end{itemize}

The general presentation of this \textit{short} documentation is:

\begin{myprestexcode}{Short explanation}{Courte explication}
\macro(*)[opt arg]{arg}

% --- [keys/clés] ------------------------------------------------------------
% key1[en] = key1[fr] = default
% key2[en] = key2[fr] = default
% ...
% ----------------------------------------------------------------------------
\end{myprestexcode}

\begin{myprestex}{Short explanation}{Courte explication}
\LaTeX
\end{myprestex}

For further details (especially regarding keys), the French documentation provides additional explanations and examples.

\smallskip

As the commands and keys are currently being anglicized, some may not be included in this documentation. Please note that the vast majority of \texttt{ProfLycee} commands are already compatible with both English and French.

\pagebreak

\tableofcontents

\pagebreak

\section{Loading, global usage}

\subsection{Loading}

\texttt{Proflycee} can be loaded with:

\begin{myprestexcode}{Loading}{Chargement}
\usepackage[options]{ProfLycee}
\usepackage[options]{ProfLycee-Light}

% --- {options} --------------------------------------------------------------
% xcolor                     := load xcolor with [table,svgnames]
% nosiunitxfr = nonsiunitxfr := do not load sinuitx with french parameters
% noamssymb   = nonamssymb   := do not load amssymb
% warningsoff = nonwarnings  := suppress warnings for unicode-math
% notikzbabel = nontikzbabel := do not load tikzbabel library
% nocancel    = noncancel    := do not load cancel
% nofa        = nonfa        := do not load fontawesome (fa5 default loaded)
% fa6                        := load fa6
% fa7                        := load fa7
% french                     := load french version (activated by default)
% english                    := load english version
% german                     := load german version
% spanish                    := load spanish version
% ----------------------------------------------------------------------------
\end{myprestexcode}

\subsection{Global usage}

Modules (libraries) can be loaded in addition to main package:

\begin{myprestexcode}{Libraries}{Librairies}
\useproflyclib{lib}

% --- {lib for full} ---------------------------------------------------------
% piton
% space       = espace
% mathwriting = ecritures
% ----------------------------------------------------------------------------

% --- {lib for light} --------------------------------------------------------
% analysis    = analyse
% listings
% probas
% stats
% arithm
% random      = aleatoire
% seq         = suites
% space       = espace
% mathwriting = ecritures
% piton
% ----------------------------------------------------------------------------
\end{myprestexcode}

Language can be loaded with:

\begin{myprestexcode}{Language}{Langue}
\setproflyclng{lng}

% --- {lng} ------------------------------------------------------------------
% fr / en / de / es
% ----------------------------------------------------------------------------
\end{myprestexcode}

\newpage

\section{Core mathematical utilities}

\subsection{Main macros}

\begin{myprestexcode}{Core calculus}{Calculs de base}
\pflnum(*)[prec]{expr}
\pflnumfrac[format]{expr}

% --- (starred) --------------------------------------------------------------
% * for real expr
% 
% --- [format] ---------------------------------------------------------------
% for math output (d/n/t)
% ----------------------------------------------------------------------------
\end{myprestexcode}

\begin{myprestex}{Core calculus}{Calculs de base}
\pflnum{2^12+1} | 
\pflnum*{4096/(12+1)} | 
\pflnum*[4]{4096/(12+1)} | 
\pflnum*[5]{sqrt(1+exp(9.75))} | 
\pflnumfrac{74+1/5} \pflnumfrac[d]{74+1/5}
\end{myprestex}

\begin{myprestexcode}{Round}{Arrondi}
\pflround(*)[prec]{expr}
\pflarrondi(*)[prec]{expr}

% --- (starred) --------------------------------------------------------------
% * for display sign
% ----------------------------------------------------------------------------
\end{myprestexcode}

\begin{myprestex}{Round}{Arrondi}
$1+\dfrac{7}{11} \approx \pflround{1+7/11} \approx \pflround[5]{1+7/11} \approx
\pflround*[2]{1+7/11}$

$\ln\big(1+\e^{4}\big) \approx \pflround[6]{log(1+exp(4))}$
\end{myprestex}

\subsection{Simple conversion between basis}

\begin{myprestexcode}{Conversion between basis}{Conversion entre bases}
\pflbasetobase{init basis}{final basis}{number}
\pflbasetobase{base départ}{base arrivée}{nombre}
\end{myprestexcode}

\begin{myprestex}{Conversion between basis}{Conversion entre bases}
\pflbasetobase{10}{2}{27} |         % 10 -> 2
\pflbasetobase{10}{16}{456} |       % 10 -> 16
\pflbasetobase{2}{16}{10011111} |   % 2  -> 16
\pflbasetobase{16}{7}{AB5F}         % 16 -> 7
\end{myprestex}

\newpage

\section{Analysis}

\subsection{Functions}

\begin{myprestexcode}{Determine Maximum}{Détermine Max}
\pfldetmax[round]{fct}{init}{step}{final}[\tmpmax][\tmpmaxvalx]
\DetermineMax[arrondi]{fct}{init}{pas}{final}[\tmpmax][\tmpmaxvalx]
\end{myprestexcode}

\begin{myprestex}{Determine Maximum}{Détermine Max}
\pfldetmax{80*x*exp(-0.2*x)}{3}{10}\tmpmax\ in \tmpmaxvalx
\end{myprestex}

\begin{myprestexcode}{Determine Minimum}{Détermine Min}
\pfldetmin[round]{fct}{init}{step}{final}[\tmpmax][\tmpmaxvalx]
\DetermineMin[arrondi]{fct}{init}{pas}{final}[\tmpmax][\tmpmaxvalx]
\end{myprestexcode}

\begin{myprestex}{Determine Minimum}{Détermine Min}
\pfldetmin[0.001]%
  {log(exp(-x)+1)+0.25*x}{1}{1.2}%
  [\mymin]%
  [\myminx]%
\mymin\ in \myminx
\end{myprestex}

\begin{myprestexcode}{Approx sol f(x)=k}{Résol approch f(x)=k}
\pflapproxsol[<keys>]{f(x)=k}[macro]
\ResolutionApprochee[<clés>]{f(x)=k}[macro]
\pflresolapproch[<clés>]{f(x)=k}[macro]

% --- [keys/clés] ------------------------------------------------------------
% prec     = Precision  = 2
% var      = Variable   = x
% interval = Intervalle = {0:10}
%
% --- [macro] ----------------------------------------------------------------
% \macrod := round by default
% \macroe := round by excess
% \macroa := round
% ----------------------------------------------------------------------------
\end{myprestexcode}

\begin{myprestex}{Approx sol f(x)=k}{Résol approch f(x)=k}
\pflapproxsol[var=t,interval=-1:2,prec=4]{3*t*exp(-0.5*t+1)=5.365}[SolA]%
$t_1 \approx \num[minimum-decimal-digits=4]{\SolAd}$\\
$t_1 \approx \num[minimum-decimal-digits=4]{\SolAe}$\\
$t_1 \approx \num[minimum-decimal-digits=4]{\SolAa}$
\end{myprestex}

\begin{myprestexcode}{Sol by ITV}{Sol par TVI}
\pflivtsol[<clés>]{f(x)}{k}
\SolutionTVI[<clés>]{f(x)}{k}
\pflsoltvi[<clés>]{f(x)}{k}

% --- [keys/clés] ------------------------------------------------------------
% fct name   = NomFct       = f
% prec       = Precision    = 2
% sol name   = NomSol       = \alpha
% stretch    = Stretch      = 1.15
% sweep      = Balayage     = false
% calculator = Calculatrice = false
% capital    = Majuscule    = true
% ----------------------------------------------------------------------------
\end{myprestexcode}

\begin{myprestex}{Sol by ITV}{Sol by TVI}
\pflivtsol[va=1.414,vb=1.415,prec=5]{x**2-2}{0}
\end{myprestex}

\begin{myprestexcode}{Canonical form}{Forme canonique}
\pflcanonicform(*)[format]{a}{b}{c}
\pflformcanoniq(*)[format]{a}{b}{c}
\FormeCanonique(*)[format]{a}{b}{c}

% --- (starred) --------------------------------------------------------------
% * for minus sign
%
% --- [format] ---------------------------------------------------------------
% for math output (d/n/t)
% ----------------------------------------------------------------------------
\end{myprestexcode}

\begin{myprestex}{Canonical form}{Forme canonique}
$2x^2+4x-6 = \pflcanonicform{2}{4}{-6} = \pflcanonicform*{2}{4}{-6}$
\end{myprestex}

\begin{myprestexcode}{Homographic form}{Forme homographique}
\pflhomogfct(*)[format]{a}{b}{c}
\pflfcthomogr(*)[format]{a}{b}{c}
\FonctionHomographique(*)[format]{a}{b}{c}

% --- (starred) --------------------------------------------------------------
% * for minus sign
%
% --- [format] ---------------------------------------------------------------
% for math output (d/n/t)
% ----------------------------------------------------------------------------
\end{myprestexcode}

\begin{myprestex}{Homographic form}{Forme homographique}
$\dfrac{11}{3x+6} =
\pflhomogfct{0}{11}{3}{6} = \pflhomogfct*{0}{11}{3}{6}$
\end{myprestex}

\subsection{Sequences}

\begin{myprestexcode}{Calc recurrence}{Calcul terme récurrence}
\pflcalcrecurr[<keys>]{fct}
\CalculTermeRecurrence[<clés>]{fct}

% --- [keys/clés] ------------------------------------------------------------
% first n  = No
% first un = UNo
% prec     = Precision  = 3
% final n = N
% ----------------------------------------------------------------------------
\end{myprestexcode}

\begin{myprestex}{Calc recurrence}{Calcul terme récurrence}
\pflcalcrecurr[prec=6,first n=0,first un=50,final n=20]{1/(x+2)}
\end{myprestex}

\begin{myprestexcode}{Threshold}{Seuil}
\pflsolthreshold[<keys>]{fct}{val}
\SolutionSeuil[<clés>]{fct}{val}
\pflsolseuil[<clés>]{fct}{val}

% --- [keys/clés] ------------------------------------------------------------
% prec       = Precision    = 3
% stretch    = Stretch      = 1.15
% sweeping   = Balayage     = false
% calc       = Calculatrice = false
% capital    = Majuscule    = true
% sign       = Sens         = {>}
% exact      = Exact        = false
% exactB     = ExactB       = false
% conclusion = Conclusion   = true
% simple     = Simple       = false
% ----------------------------------------------------------------------------
\end{myprestexcode}

\begin{myprestex}{Threshold}{Seuil}
\pflsolthreshold[prec=4,first n=1,first un=2]{1+(1+x**2)/(1+x)}{10}
\end{myprestex}

\subsection{Integrals}

\begin{myprestexcode}{Integral approx}{Approch intégrale}
\pflapproxintegral[<keys>]{fct}{a}{b}
\IntegraleApprochee[<clés>]{fct}{a}{b}
\pflintegrapproch[<clés>]{fct}{a}{b}

% --- [keys/clés] ------------------------------------------------------------
% prec         = Precision    = 3
% nb subdiv    = NbSubDiv     = 10
% method       = Methode      = Simpson
% raw result   = ResultatBrut = false
% disp formula = AffFormule   = false
% sign         = Signe        = {\approx}
% expr         = Expr         = f(x)
% var          = Variable     = x
% ----------------------------------------------------------------------------
\end{myprestexcode}

\begin{myprestex}{Integral approx}{Approch intégrale}
\pflapproxintegral[raw result]{sqrt(x)}{4}{10}
\end{myprestex}

\begin{myprestexcode}{Raw value of integral}{Valeur 'brute' d'une intégrale}
\pflcalcintegral(*)[prec]<num of subdiv>{fct}{a}{b}
\CalcIntegrale(*)[prec]<num of subdiv>{fct}{a}{b}
\pflvalintegr(*)[prec]<nb de subdiv>{fct}{a}{b}

% --- (starred) --------------------------------------------------------------
% * for siunitx output
% ----------------------------------------------------------------------------
\end{myprestexcode}

\begin{myprestex}{Raw value of integral}{Valeur 'brute' d'une intégrale}
\CalcIntegrale{80*x*exp(-0.2*x)}{1}{20}
\end{myprestex}

\begin{myprestexcode}{Average value}{Valeur moyenne}
\pflaveragevalintegral(*)[prec]<num of subdiv>{fct}{a}{b}
\pflvalmoyintegr(*)[prec]<nb de subdiv>{fct}{a}{b}

% --- (starred) --------------------------------------------------------------
% * for siunitx output
% ----------------------------------------------------------------------------
\end{myprestexcode}

\subsection{Signs, convexity}

\begin{myprestexcode}{Sign diagram}{Schémas signes}
\pflminisignplot*[<keys>]<tikz options>
\pflschemasignes*[<clés>]<options tikz>
\MiniSchemaSignes*[<clés>]<options tikz>

% --- [keys/clés] ------------------------------------------------------------
% code   = Code    = da+
% color  = Couleur = red
% width  = Largeur = 2
% height = Hauteur = 1
% border = Cadre   = true
% roots  = Racines = 2
% ----------------------------------------------------------------------------
\end{myprestexcode}

\begin{myprestex}{Sign diagram}{Schémas signes}
\pflminisignplot*[code=pa+d+,roots={1/2},color=orange]
\end{myprestex}

\begin{myprestexcode}{Sign diagram with tkztab}{Schémas signes avec tkztab}
\pflminisignplottkztab[<keys>]{line}[scale][hoffset]
\pflschemasignestkztab[<clés>]{line}[scale][hoffset]
\MiniSchemaSignesTkzTab[<clés>]{ligne}[echelle][décal H]

% --- [keys/clés] ------------------------------------------------------------
% code   = Code    = da+
% color  = Couleur = red
% width  = Largeur = 2
% height = Hauteur = 1
% border = Cadre   = true
% roots  = Racines = 2
% ----------------------------------------------------------------------------
\end{myprestexcode}

\begin{myprestex}{Sign diagram with tkztab}{Schémas signes avec tkztab}
\begin{tikzpicture}
  \tkzTabInit[espcl=2]%
    {$x$/1,$-2x+5$/1,$2x+4$/1,$p(x)$/1}{$-\infty$,$-2$,${2.5}$,$+\infty$}
  \tkzTabLine{,+,t,+,z,-,}
  \tkzTabLine{,-,z,+,t,+,}
  \tkzTabLine{,-,z,+,z,-,}
  \pflminisignplottkztab[code=da-,roots={\tfrac{5}{2}},color=blue]{1}
  \pflminisignplottkztab[code=da+,roots={-2},color=purple]{2}
  \pflminisignplottkztab[code=pa-d+,roots={-2/{\tfrac{5}{2}}},color=orange]%
    {3}[0.85][2]
\end{tikzpicture}
\end{myprestex}

\begin{myprestexcode}{Convexity line with tkztab}{Ligne de convexité avec tkztab}
\pfltkztabcvx(*){line}{'cells'}
\tkzTabLineConvex(*){ligne}{liste 'cases'}

% --- (starred) --------------------------------------------------------------
% * for text version
% ----------------------------------------------------------------------------
\end{myprestexcode}

\begin{myprestex}{Convexity line with tkztab}{Ligne de convexité avec tkztab}
\begin{tikzpicture}
  \tkzTabInit[lgt=3,espcl=2.5]%
    {$x$/1,$f''(x)$/1,Convexity\\of $f$/2}{$0$,$1$,$3$,$5$}
  \tkzTabLine{,+,z,-,z,+,}
  \pfltkztabcvx{2}{,cvx,i*,ccv,i*,cvx,}
\end{tikzpicture}
\end{myprestex}

\newpage

\section{Code presentation}

\subsection{Python Code}

\begin{myprestexcode}{Environment python code listings}{Environnement code python listings}
\begin{pflpythoncodelst}*[<keys>]{tcbox options} ... \end{pflpythoncodelst}
\begin{CodePythonLst}*[<clés>]{options tcbox}   ...  \end{CodePythonLst}

% --- (starred) --------------------------------------------------------------
% * for removing line numbers
%
% --- [keys/clés] ------------------------------------------------------------
% width      = Largeur    = \linewidth
% first line = PremLigne  = 1
% hsep       = EspaceNum  = 14pt
% ----------------------------------------------------------------------------
\end{myprestexcode}

\begin{myprestex}{Environment python code listings}{Environnement code python listings}
\begin{pflpythoncodelst}*[width=0.75\linewidth]{flush right}
nb = int(input("n = "))
if (nb %7 == 0) :
	print(f"{nb} is ok")
#endif

def f(x) :
	return x**2
\end{pflpythoncodelst}
\end{myprestex}

\begin{myprestexcode}{Environment python code piton}{Environnement code python piton}
\begin{pflpitoncode}[<keys>]{tcbox options} ... \end{pflpitoncode}
\begin{CodePiton}[<clés>]{options tcbox}    ... \end{CodePiton}
\pflpitoncodefile[<keys>]{tcbox options}
\CodePitonFichier[<clés>]{options tcbox}

% --- [keys/clés] ------------------------------------------------------------
% width        = Largeur        = \linewidth
% align        = Alignement     = center
% lines        = Lignes         = true
% gobble       = Gobble         = {}
% font size    = TaillePolice   = \footnotesize
% watermark    = Filigrane      = false
% style        = Style          = classic/modern/(light)carbon/tab
% border       = Cadre          = true
% title bar    = BarreTitre     = true
% title        = Titre          = ...
% vsep         = EspacementV    = 0.5\baselineskip
% num color    = CouleurNombres = {orange!75!black}
% color        = Couleur        = red
% carbon icons = IconesCarbon   = ...
% ----------------------------------------------------------------------------
\end{myprestexcode}

\begin{myprestex}{Environment python code piton}{Environnement code python piton}
\begin{pflpitoncode}%
	[gobble=tabs,width=12cm,style=lightcarbon]{}
def f(x) :
	return x**2
\end{pflpitoncode}
\end{myprestex}

\begin{myprestexcode}{Console python piton}{Console python piton}
\begin{pflpitonconsole}[<keys>]{tcbox options} ... \end{pflpitonconsole}
\begin{ConsolePiton}[<clés>]{options tcbox}    ... \end{ConsolePiton}

% --- [keys/clés] ------------------------------------------------------------
% logo            = Logo            = true
% width           = Largeur         = \linewidth
% align           = Alignement      = {flush left}
% num color       = CouleurNombres  = {orange!75!black}
% beforeafter sep = CoeffAvantApres = 0.75
% ----------------------------------------------------------------------------
\end{myprestexcode}

\begin{myprestex}{Console python piton}{Console python piton}
\begin{pflpitonconsole}{}
1+1
2**10
\end{pflpitonconsole}
\end{myprestex}

\begin{myprestexcode}{Thonny piton}{Thonny piton}
\begin{pflpitonthonnyeditor}[<keys>]{...}...\end{pflpitonthonnyeditor}
\begin{PitonThonnyEditor}[<clés>]{...}...\end{PitonThonnyEditor}
\pflpitonthonnyeditorfile[<keys>]{...}
\PitonThonnyEditorFichier[<clés>]{...}

% --- [keys/clés] ------------------------------------------------------------
% filename     = NomFichier     = {script.py}
% consolename  = NomConsole     = console
% gobble       = Gobble         = {}
% num color    = CouleurNombres = {orange!75!black}
% width        = Largeur        = \linewidth
% ----------------------------------------------------------------------------
\end{myprestexcode}

\begin{myprestex}{Thonny piton}{Thonny piton}
\begin{pflpitonthonnyeditor}[NomFichier=tpcapytale.py,Largeur=12cm]{}
nb = int(input("n = "))
if (nb %7 == 0) :
	print(f"{nb} is ok")
#endif

def f(x) :
	return x**2
\end{pflpitonthonnyeditor}
\end{myprestex}

\begin{myprestexcode}{Thonny console piton}{Thonny console piton}
\begin{pflpitonthonnyconsole}[<keys>]{tcbox} ... \end{pflpitonthonnyconsole}
\begin{PitonThonnyConsole}[<clés>]{tcbox} ...    \end{PitonThonnyConsole}

% --- [keys/clés] ------------------------------------------------------------
% filename      = NomFichier      = {script.py}
% consolename   = NomConsole      = console
% console intro = IntroConsole    = {Python 3.11.6 /usr/bin/python}
% gobble        = Gobble          = {}
% num color     = CouleurNombres  = {orange!75!black}
% width         = Largeur         = \linewidth
% ----------------------------------------------------------------------------
\end{myprestexcode}

\begin{myprestex}{Thonny console piton}{Thonny console piton}
\begin{pflpitonthonnyconsole}[consolename={python 3.8.10},width=12cm]{}
#Run script.py
from random import randint
nb1 = randint(1,100000)
nb2= randint(1,100000)
print(nb1, nb2)
\end{pflpitonthonnyconsole}
\end{myprestex}

\subsection{Pseudo-Code}

\begin{myprestexcode}{Environment pseudocode piton}{Environnement pseudocode piton}
\begin{pflpseudocodepiton}[<keys>]{tcbox} ... \end{pflpseudocodepiton}
\begin{PseudoCodePiton}[<clés>]{tcbox}    ... \end{PseudoCodePiton}
\pflpitonpseudocodefile[<keys>]{tcbox}
\PseudoCodePitonFichier[<clés>]{tcbox}

% --- [keys/clés] ------------------------------------------------------------
% width       = Largeur      = \linewidth
% align       = Alignement   = center
% lines       = Lignes       = true
% gobble      = Gobble       = {}
% font size   = TaillePolice = \footnotesize
% watermark   = Filigrane    = false
% style       = Style        = Classique
% border      = Cadre        = true
% title bar   = BarreTitre   = true
% title       = Titre        = ...
% vsep        = EspacementV  = 0.5\baselineskip
% color       = Couleurs     = true
% ----------------------------------------------------------------------------
\end{myprestexcode}

\begin{myprestex}{Environment pseudocode piton}{Environnement pseudocode piton}
%in french...
\begin{PseudoCodePiton}[Lignes=false,Filigrane=false,Couleurs=true]{}
Algorithme : Périmètre de rectangles
Variables  : Long, Larg, Perim (réels)
             Choix (chaîne) # pour refaire ou non l'exécution
Début
	# initialisation de l'indicateur pour entrer dans la boucle
	Choix ← "o"
	# boucle TantQue permettant de refaire le traitement selon le choix
	TantQue Choix = "o" Faire
		# traitement
		Afficher("Donner les dimensions du rectangle") et Saisir(Long,Larg)
		Perim ← 2 × (Long + Larg)
		Afficher("Le périmètre du rectangle est", Perim)
		#saisie du choix de recommencer ou non
	Afficher("Voulez-vous continuer (o/n) ?") et Saisir(Choix)
	FinTantQue
Fin
\end{PseudoCodePiton}
\end{myprestex}

%\newpage
%
%\section{Geometry}

\newpage

\section{Statistics, probabilities}

\subsection{Statistics}

\begin{myprestexcode}{Stats parameters}{Paramètres stats}
\DeterminerParamStats(*){liste}[\macromin]...[\macromax]
\pflparamstats(*){liste}[\macromin]...[\macromax]

% --- (starred) --------------------------------------------------------------
% * for data with regroupment
% ----------------------------------------------------------------------------
\end{myprestexcode}

\begin{myprestex}{Stats parameters}{Paramètres stats}
\def\listedonnees%
{0.7426,0.7429,0.7435,0.7448,0.7456,0.7478,%
	0.7478,0.748,0.7487,0.7494,0.7494,0.7496}
\pflparamstats{\listedonnees}
\monmin/\monquartileun/\mamediane/\monquartiletrois/\monmax
\end{myprestex}

\subsection{Probabilities}

\begin{myprestexcode}{Proba tree}{Arbre de probas}
\pflprobtree[<keys>]{data}
\ArbreProbasTikz[<clés>]{donnees}
\pflarbreprobas[<clés>]{donnees}

% --- [keys/clés] ------------------------------------------------------------
% unit           = Unite          = 1cm
% level sep      = EspaceNiveau   = 3.25
% child sep      = EspaceFeuille  = 1
% type           = Type           = 2x2
% font           = Police         = ...
% prob font      = PoliceProbas   = ...
% sloped         = InclineProbas  = true
% arrows         = Fleche         = false
% arc style      = StyleTrait     = {}
% arc thickness  = EpaisseurTrait = semithick
% prob pos       = PositionProbas = {}
% bg color       = CouleurFond    = white
% arc color      = CouleurTraits  = black
% prob color     = CouleurProbas  = black
% node color     = CouleurNoeuds  = black
% root           = Racine         = {}
% ----------------------------------------------------------------------------
\end{myprestexcode}

\begin{myprestex}{Proba tree}{Arbre de probas}
\def\ArbreDeuxDeux{
  $A$/\num{0.5}/,
  $B$/\num{0.4}/,
  $\overline{B}$/.../,
  $\overline{A}$/.../,
  $B$/.../,
  $\overline{B}$/$\frac{1}{3}$/
}
\pflprobtree{\ArbreDeuxDeux}
\end{myprestex}
\subsection{Random}

\begin{myprestexcode}{Randint}{Randint}
\pflrandnum[n]{a}{b}{macro}
\NbAlea[n]{a}{b}{macro}
\pflnbalea[n]{a}{b}{macro}
\end{myprestexcode}

\begin{myprestexcode}{Rand number}{Nombre aléatoire}
\pflvarrandnum{macro}{calculs}
\VarNbAlea{macro}{calculs}
\pflvarnbalea{macro}{calculs}
\end{myprestexcode}

\begin{myprestex}{Randint/randnum}{Randint/randum}
Integer between 1 and 50: \pflrandnum{1}{50}{\FirstNbAlea}
\FirstNbAlea

\pflvarrandnum{\SecondNbAlea}{\FirstNbAlea+randint(0,10)}
\SecondNbAlea

Float between 0 and 10.999... : \pflrandnum[3]{0}{10}{\FirstFloatAlea}
\FirstFloatAlea
\end{myprestex}

\begin{myprestexcode}{Randint list}{Liste d'entiers aléatoires}
\pflrandintnumbers[<keys>]{macro}
\TirageAleatoireEntiers[<clés>]{macro}
\pfltiragealeaent[<clés>]{macro}

% --- [keys/clés] ------------------------------------------------------------
% min value  = ValMin     = 1
% max value  = ValMax     = 50
% nb values  = NbVal      = 6
% sep        = Sep        = {,}
% sort       = Tri        = non
% repeat     = Repetition = false
% ----------------------------------------------------------------------------
\end{myprestexcode}

\begin{myprestex}{Randint list}{Liste d'entiers aléatoires}
\pflrandintnumbers{\FirstListAlea}\FirstListAlea
\end{myprestex}

\begin{myprestexcode}{Choice tree}{Arbre "sans remise"}
\pflchoicetree[<keys>]<tikz options>{list}
\ArbreChoix[clés]<options tikz>{liste}
\pflarbrechoix[<clés>]<options tikz>{liste}

% --- [keys/clés] ------------------------------------------------------------
% level sep    = EspaceNiveaux   = 2.25
% child sep    = EspaceFeuilles  = 0.5
% scale        = Echelle         = 1
% repeat       = Repet           = {}
% help         = Notice          = {}
% help rules   = TraitsNotice    = false
% level color  = CouleursNiveaux = black
% show results = AffResultats    = false
% result sep   = SepResultats    = {}
% ----------------------------------------------------------------------------
\end{myprestexcode}

\begin{myprestex}{Choice tree}{Arbre "sans remise"}
\pflchoicetree{C/S,P/M/G,N/B}
\end{myprestex}

\begin{myprestexcode}{Tree no repeating}{Arbre "sans remise"}
\pflchoicetreenorepl<tikz options>{list}
\ArbreChoixSansRemise[clés]<options tikz>{liste}
\pflarbrechoixssremise[<clés>]<options tikz>{liste}

% --- [keys/clés] ------------------------------------------------------------
% level sep    = EspaceNiveaux   = 2.25
% child sep    = EspaceFeuilles  = 0.5
% scale        = Echelle         = 1
% repeat       = Repet           = {}
% help         = Notice          = {}
% help rules   = TraitsNotice    = false
% level color  = CouleursNiveaux = black
% show results = AffResultats    = false
% result sep   = SepResultats    = {}
% ----------------------------------------------------------------------------
\end{myprestexcode}

\begin{myprestex}{Tree no repeating}{Arbre "sans remise"}
\pflchoicetreenorepl[level color={gray,olive,red,blue}]{1,3,5,7}
\end{myprestex}

\newpage

\section{Analytical geometry}

\subsection{Cartesian equation of a plane}

\begin{myprestexcode}{Cartesian equation of plane}{Équation cartésienne d'un plan}
\pflcartplaneq[keys](normal vect)(point)
%or
\pflcartplaneq[keys](dir vect 1)(dir vect 2)(point)
%or
\pflcartplaneq[keys](point1)(point2)(point3)

\TrouveEqCartPlan[keys]......

% --- [keys/clés] ------------------------------------------------------------
% coeff fmt    = OptionCoeffs  = d
% simpl coeffs = SimplifCoeffs = false
% factor       = Factor        = true
% ----------------------------------------------------------------------------
\end{myprestexcode}

\begin{myprestex}{Cartesian equation of plane}{Équation cartésienne d'un plan}
$\pflcartplaneq(1;2;3)(4,5,6)$\\
$\pflcartplaneq(1/2;2/3;3/5)(4,5,6)$ or
$\pflcartplaneq[simpl coeffs](1/2;2/3;3/5)(4,5,6)$\\
$\pflcartplaneq(2,0,1)(3,1,1)(1,-2,0)$
\end{myprestex}

\subsection{Cartesian equation of a sphere}

\begin{myprestexcode}{Cartesian equation of sphere}{Équation cartésienne d'une sphère}
\pfleqsphere[keys](point)(radius or diam point or cart plan)
\TrouveEqSphere[clés](point)(rayon ou diamètre ou éq cart)

% --- [keys/clés] ------------------------------------------------------------
% coeff fmt    = OptionCoeffs  = d
% simpl coeffs = SimplifCoeffs = false
% diameter     = Diam          = false
% expand       = Develop       = true
% sort         = Sort          = true
% ----------------------------------------------------------------------------
\end{myprestexcode}

\begin{myprestex}{Cartesian equation of sphere}{Équation cartésienne d'une sphère}
%center + radius
$\pfleqsphere(1,1,1)(2)$\\
$\pfleqsphere[expand=false](1,1,1)(2)$\\
$\pfleqsphere[sort=false](1,1,1)(2)$
\end{myprestex}

\begin{myprestex}{Cartesian equation of sphere}{Équation cartésienne d'une sphère}
%center + point
$\pfleqsphere(3,1,4)(0,0,0)$\\
$\pfleqsphere[expand=false](3,1,4)(0,0,0)$\\
$\pfleqsphere[sort=false](3,1,4)(0,0,0)$
\end{myprestex}

\begin{myprestex}{Cartesian equation of sphere}{Équation cartésienne d'une sphère}
%diameter
$\pfleqsphere[diameter](3,1,4)(0,0,0)$\\
$\pfleqsphere[diameter,expand=false](3,1,4)(0,0,0)$\\
$\pfleqsphere[diameter,sort=false](3,1,4)(0,0,0)$
\end{myprestex}

\begin{myprestex}{Cartesian equation of sphere}{Équation cartésienne d'une sphère}
%center + tangent plan
$\pfleqsphere[](1,2,4)(2x+3y-3z+1=0)$\\
$\pfleqsphere[expand=false](1,2,4)(2x+3y-3z+1=0)$\\
$\pfleqsphere[sort=false](1,2,4)(2x+3y-3z+1=0)$
\end{myprestex}

\newpage

\section{Arithmetic}

\subsection{Conversions}

\begin{myprestexcode}{Conversion to dec}{Conversion en b10}
\pflconvtodec[keys]{number}
\pflconvversdec[clés]{nb}
\ConversionVersDec[clés]{nb}

% --- [keys/clés] ------------------------------------------------------------
% base from  = BaseDep  = 2
% show base  = AffBase  = true
% details    = Details  = true
% zeros      = Zeros    = true
% ----------------------------------------------------------------------------
\end{myprestexcode}

\begin{myprestex}{Conversion to dec}{Conversion en b10}
$\ConversionVersDec[zeros=false,base from=16]{AC0DC}$
\end{myprestex}

\begin{myprestexcode}{Conversion from dec}{Conversion depuis b10}
\pflconvdepuisdix[clés]{nombre}{base arrivée}
\pflconvfromten[keys]{number}{final basis}
\ConversionDepuisBaseDix[clés]{nombre}{base arrivée}

% --- [keys/clés] ------------------------------------------------------------
% color     = Couleur    = red
% hsep      = DecalH     = 2pt
% vsep      = DecalV     = 3pt
% rect      = Rect       = true
% node      = Noeud      = EEE
% res color = CouleurRes = false
% ----------------------------------------------------------------------------
\end{myprestexcode}

\begin{myprestex}{Conversion from dec}{Conversion depuis b10}
\ConversionDepuisBaseDix{78}{2}
\end{myprestex}

\subsection{Prime numbers, PGCD}

\begin{myprestexcode}{Isprime ?}{Est premier ?}
\pflisprimebool{expr}
\pflboolestpremier{expr}
\end{myprestexcode}

\begin{myprestex}{Isprime ?}{Est premier ?}
\pflisprimebool{25457} and \pflisprimebool{25458}
\end{myprestex}

\begin{myprestexcode}{Number factor}{Décompo facteurs premiers}
\pflnumberfact(*){expr}
\pfldecompnb(*){expr}

% --- (starred) --------------------------------------------------------------
% * for 1 power
% ----------------------------------------------------------------------------
\end{myprestexcode}

\begin{myprestex}{Number factor}{Décompo facteurs premiers}
$24458=\pflnumberfact{24458}=\pflnumberfact*{24458}$
\end{myprestex}

\begin{myprestexcode}{Gcd presentation}{Présentation PGCD}
\pflprespgcd[keys]{a}{b}
\PresentationPGCD[clés]{a}{b}

% --- [keys/clés] ------------------------------------------------------------
% color      = Couleur            = red
% rect sep   = DecalRect          = 2pt
% rect       = Rectangle          = true
% node       = Noeud              = FFF
% res color  = CouleurResultat    = false
% disp ccl   = AfficheConclusion  = true
% disp delim = AfficheDelimiteurs = true
% ----------------------------------------------------------------------------
\end{myprestexcode}

\begin{myprestex}{Gcd presentation}{Présentation PGCD}
\pflprespgcd[res color,color=CouleurVertForet]{1250}{450}.
\end{myprestex}

\subsection{Divisors}

\begin{myprestexcode}{Tree divisors}{Arbre des diviseurs}
\pfltreediv[keys]<tikz options>{expr}
\pflarbrediv[clés]<options tikz>{expr}
\ArbreDiviseurs[clés]<options tikz>{expr}

% --- [keys/clés] ------------------------------------------------------------
% level sep     = EspaceNiveau   = 2.25
% child sep     = EspaceFeuille  = 0.66
% details       = Details        = true
% details color = CouleurDetails = red
% scale         = Echelle        = 1
% arrows        = Fleches        = true
% ----------------------------------------------------------------------------
\end{myprestexcode}

\begin{myprestex}{Tree divisors}{Arbre des diviseurs}
\pfltreediv{60}
\end{myprestex}

\subsection{Prime decomposition}

\begin{myprestexcode}{Prime decomposition}{Décomposition d'un entier}
\pflfactornb[keys]{number}
\DecompFactPremiers[clés]{nombre}

% --- [keys/clés] ------------------------------------------------------------
% long      = Longue         = false
% onepower  = PuissanceUn    = false
% ----------------------------------------------------------------------------
\end{myprestexcode}

\begin{myprestex}{Prime decomposition}{Décomposition d'un entier}
\pflfactornb{2781240}\\
\pflfactornb[onepower]{2781240}\\
\pflfactornb[long]{2781240}
\end{myprestex}

\begin{myprestexcode}{Prime presentation}{Présentation de la décomposition}
\pflprimefactpres[keys]{number}
\pflpresfactprem[clés]{nombre}
\PresFactPremiers[clés]{nombre}

% --- [keys/clés] ------------------------------------------------------------
% rulecolor = CouleurTrait   = black
% rulethick = EpaisseurTrait = 0.4pt
% endbox    = EncadreFin     = false
% ----------------------------------------------------------------------------
\end{myprestexcode}

\begin{myprestex}{Prime presentation}{Présentation de la décomposition}
\pflprimefactpres{2781240}\\
\pflprimefactpres[endbox,rulecolor=red,rulethick=0.8pt]{2781240}
\end{myprestex}

\subsection{Miscellaneous}

\begin{myprestexcode}{Period fraction}{Fraction pédiodique}
\pflperiodfraction[keys]{before period}{period}
\pflfracperiod[clés]{avant période}{période}
\FractionPeriode[clés]{avant période}{période}

% --- [keys/clés] ------------------------------------------------------------
% statement = Enonce   = true
% d                    = true
% unknown   = Inconnue = x
% sol       = Solution = false
% raw       = Brut     = true
% simple    = Simple   = false
% ----------------------------------------------------------------------------
\end{myprestexcode}

\begin{myprestex}{Period fraction}{Fraction pédiodique}
\pflperiodfraction[sol]{45.1}{23}
\end{myprestex}

\newpage

\section{Recreations}

\subsection{Fractals}

\begin{myprestexcode}{Fractals, with \TikZ}{Fractales, avec \TikZ}
\pfltkzfractal[type=...,<keys>]<tikZ options>
\FractaleTikz[Type=...,<clés>]<options tikz>
\pflfractaltikz[Type=.,<clés>]<options tikz>

% --- [type] -----------------------------------------------------------------
% type        = Type         = Koch or Sierp
%
% --- [keys/clés] ------------------------------------------------------------
% thickness   = Epaisseur    = 0.6pt
% color       = Couleur      = black
% side width  = LongueurCote = 3
% step        = Etape        = 1
% fill        = Remplir      = false
% fill style  = Remplissage  = lightgray
% start       = Depart       = {(0,0)}
% v align     = AlignV       = false
% offset      = Offset       = 2pt
% ----------------------------------------------------------------------------
\end{myprestexcode}

\begin{myprestex}{Fractals, with \TikZ}{Fractales, avec \TikZ}
\pfltkzfractal[type=Koch]
\pfltkzfractal[type=Sierp]
\end{myprestex}

\subsection{House of cards}

\begin{myprestexcode}{House of cards}{Château de cartes}
\pflhouseofcards[keys]{levels}<tikz options>
\pflchateaucartes[clés]{étages}<options tikz>
\ChateauCartes[clés]{étages}<options tikz>

% --- [keys/clés] ------------------------------------------------------------
% scale       = Echelle       = 1
% deco color  = CouleurDeco   = black
% rounded     = Arrondi       = true
% y angle     = AngleY        = 20
% x angle     = AngleX        = 8
% bottom      = Bas           = false
% legend      = Legende       = false
% legend font = PoliceLegende = \normalsize\normalfont
% deco        = Deco          = remplir
% ----------------------------------------------------------------------------
\end{myprestexcode}

\begin{myprestex}{House of cards}{Château de cartes}
\pflhouseofcards[deco color=teal,deco=hatch,scale=2,legend]{4}
\end{myprestex}

\subsection{Stack of balls}

\begin{myprestexcode}{Stack of balls}{Empilement de balles}
\pflballstacking[keys]{levels}<tikz options>
\pflempilballes[clés]{niveaux}<options tikz>
\EmpilementBalles[clés]{niveaux}<options tikz>

% --- [keys/clés] ------------------------------------------------------------
% color    = Couleur  = gray
% rotation = Rotation = -5
% scale    = Echelle  = 1
% ----------------------------------------------------------------------------
\end{myprestexcode}

\begin{myprestex}{Stack of balls}{Empilement de balles}
\pflballstacking[scale=0.33,color=red]{3-8}
\end{myprestex}

%\newpage

%\section{Index}

%\printindex
\end{document}